%% P31 --- Mutual information
%%
%% WHAT THE NEIGHBOURS SPENT.  `grep -rn 'prog:P31' programs/en/*.tex` prints
%% the list with its contexts and it was the first thing run.  P30 built the
%% whole of the machinery -- mutual information IS its divergence between a
%% joint and the product of the margins -- and its closing frame says where
%% the length goes: "P31 defines it, and spends most of its length on WHY THE
%% CLAIMS PEOPLE MAKE WITH IT ARE USUALLY NOT SUPPORTED."  Its header hands
%% over the data-processing inequality.  P24 built the arrival and wrote the
%% pointer: its Z/W pair has covariance exactly zero and "P31 builds the
%% quantity that does see it".  P29 owns entropy; P23 owns independence; P26
%% owns bias; P27 owns what a finite sample does to a comparison.
%%
%% SO SECTIONS 1 TO 3 ARE SHORT BY CONSTRUCTION and the length belongs in 4
%% and 5, which is what P30's closing frame already said.
%%
%% NOT SPENT HERE: P32 checks whether the independence hypotheses in this
%% book's derivations survive training, on an assembled architecture.
%%
%% Twin: programs/pl/P31-mutual-information.tex.  Frame for frame, macro for
%% macro: parity.py compares the two in order, so a sentence rebuilt in Polish
%% has to carry the same maths spans and the same numeric literals in the same
%% sequence.  A digit stays a digit and a word stays a word.

\program{Mutual information}
\label{prog:P31}
\index{mutual information}

\begin{outcomes}
\outcome{1--6}{say what a table of pairwise correlations cannot settle, and
  name the quantity that settles it}
\outcome{7--10}{say why the quantity can never be negative, when it is zero,
  and why it is symmetric}
\outcome{11--15}{say why it is symmetric when the two conditional entropies
  it is built from are not}
\outcome{16--25}{say which way a plug-in estimate is wrong on data with no
  dependence, and how many items it takes for that to stop mattering}
\outcome{26--34}{say what post-processing can and cannot do to information,
  and read the same inequality from both of its ends}
\outcome{35--38}{say what would have to be measured before an
  interpretability claim about a layer is supported}
\end{outcomes}

\begin{quiz}
\item Two quantities have a correlation of exactly zero. What have you
  learnt about whether one determines the other? \teachesat{1--2}
  \answerto{Nothing. Program~\ref{prog:P24} builds a pair whose correlation
  is exactly zero and where one is a function of the other; this program
  builds the quantity that sees it.}
\item Mutual information is symmetric. Are the two conditional entropies it
  can be written with also symmetric? \teachesat{12--13}
  \answerto{No, and they can differ enormously \dash{} one of them can be
  exactly zero while the other is not. The same number is reached from two
  different pieces.}
\item You estimate the mutual information between two quantities that are
  genuinely independent, from a finite sample. What do you get?
  \teachesat{16--17}
  \answerto{A positive number, in expectation, at every sample size. The
  truth is zero and the estimator does not return it.}
\item Somebody improves their probe and reports a higher score. What does
  that say about the layer? \teachesat{28--29}
  \answerto{Nothing new. A probe is post-processing, so its score was
  always bounded by what the layer carried; a better probe is a statement
  about the probe.}
\item And if a probe scores badly? \teachesat{30--31}
  \answerto{Also nothing, in the other direction. The bound is a lower
  bound and it is loose by an amount nobody measures.}
\end{quiz}

% ============================================================
\section{What two margins do not say}
\label{sec:P31-margins}
\index{mutual information!and zero correlation}

\begin{fr}
Program~\ref{prog:P24} left a pair on the table and a promise attached to it.

$Z$ is $-1$, $0$ or $1$, each with probability $\frac{1}{\val{p24.tri.n}}$,
and $W = Z^{2}$. It proved, exactly over fractions, that
$\Cov(Z, W) = 0$.

It also pointed out that $W$ is a \emph{function} of $Z$, so the dependence
is total. Being told $Z$ tells you $W$ outright.

Before reading on: what does a correlation of zero, on that pair, entitle
you to conclude about whether the two are independent?
\dotline
\nextframe
\end{fr}

\begin{fr}
\ans{Nothing whatever}

A correlation measures the straight-line part of a relationship and there is
no straight-line part here: the positive and negative $Z$ pull the average of
$ZW$ in opposite directions and cancel exactly. Everything else about the
relationship is invisible to it.

That is not a defect in the correlation. It is a statement about what the
number was ever measuring \dash{} and Program~\ref{prog:P23} said the same
thing from the other side, that a table of pairwise correlations cannot
settle independence.

It is worth saying what the correlation is \emph{not}. It is not useless and
it is not broken: on a pair whose relationship really is a straight line it
measures exactly the right thing, and it measures it cheaply. What it cannot
do is answer a question it was never asked, and \enquote{are these two
independent} is that question.

So what you need is a quantity that answers the question
Program~\ref{prog:P23} \emph{did} define independence with. What was that
question?
\dotline
\nextframe
\end{fr}

\begin{fr}
\ans{Does being told one of them move the distribution of the other?}

Here is that question written as arithmetic.

If $X$ and $Y$ are independent then the joint distribution is the product of
the margins, $p(x, y) = p(x)\,p(y)$, at every pair. If they are not, it is
not.

So the whole of the dependence is the gap between two distributions on the
same outcomes: the joint you have, and the product of the margins you would
have had.

You spent the whole of Program~\ref{prog:P30} on a quantity that measures a
gap between two distributions. What is it, and what should go in its two
arguments here?
\dotline
\nextframe
\end{fr}

\begin{fr}
\ans{The divergence, with the joint first and the product of the margins
second}

\[ I(X; Y) \;=\; \KL{p(x, y)}{p(x)\,p(y)} \]

and that is the definition. It is called the \textbf{mutual information} of
$X$ and $Y$.

Nothing new has been defined. It is Program~\ref{prog:P30}'s object applied
to one particular pair of distributions, and everything that program proved
about it holds here without a further word.

\begin{note}
The order matters, as it always does with that quantity, and the joint goes
first. Section~4 is about what happens when the second argument has a gap the
first does not, so it is worth knowing now which one is which.
\end{note}
\end{fr}

\begin{fr}
Apply it to the pair Program~\ref{prog:P24} left you.

The joint puts $\frac{1}{\val{p24.tri.n}}$ on each of $(-1, 1)$, $(0, 0)$ and
$(1, 1)$ and nothing anywhere else. The margins are $Z$ uniform on three
values and $W$ on two.

Work it out, or take it on trust for one frame: the answer is
$\val{p31.zw.mi}$ nats.

Now the question worth stopping for. Where has that number come from
\dash{} what is it the uncertainty \emph{of}?
\dotline
\nextframe
\end{fr}

\begin{fr}
\ans{It is exactly $H(W)$ \dash{} all of $W$'s uncertainty}

$\val{p31.zw.mi}$ nats, which is $\val{p31.zw.mi.bits}$ bits, and it is
$H(W)$ to the last decimal the script computes.

That is not a coincidence and the script asserts the identity rather than
the figure: $W$ is a function of $Z$, so being told $Z$ leaves \emph{nothing}
of $W$ undetermined, and the amount you learn about $W$ from $Z$ is
therefore everything $W$ had.

\begin{trapbox}
\textbf{\enquote{Zero correlation, so there is nothing to find} is now two
numbers rather than an argument.}

On the same pair, in the same frame: the correlation is $0$ and the mutual
information is $\val{p31.zw.mi}$ nats, which is every nat $W$ possesses.
Program~\ref{prog:P24} said a correlation matrix full of zeros is not
evidence that a set of features carries no shared information. This is what
that sentence costs when it is ignored.
\end{trapbox}

\mermaidfig{p31-what-a-margin-drops}{What the two margins throw away}{margins alone}
\end{fr}

% ============================================================
\section{Why it cannot be negative}
\label{sec:P31-nonneg}
\index{mutual information!non-negativity}

\begin{fr}
Everything Program~\ref{prog:P30} proved is now yours for nothing, because
this is that quantity and not a new one.

It proved the divergence is never negative, by applying
Program~\ref{prog:P19}'s Jensen to the ratio, and that it is zero exactly
when the two distributions agree.

Read both of those with the joint and the product of the margins in the two
slots. What do they say about $I(X; Y)$?
\dotline
\nextframe
\end{fr}

\begin{fr}
\begin{ansblock}
$I(X; Y) \ge 0$ always, and $I(X; Y) = 0$ \textbf{exactly when the joint
equals the product of its margins} \dash{} which is
Program~\ref{prog:P23}'s definition of independence, written out.
\end{ansblock}

So the quantity is zero if and only if the two are independent, and positive
otherwise. That is the property the correlation did not have, and it took no
new mathematics to get: it is a theorem about a divergence, read on a
particular pair of arguments.

Checked rather than asserted: over all $\val{p31.sym.trials}$ joint
distributions on a rational grid, not one comes out negative.
\end{fr}

\begin{fr}
There is one more property, and it is worth putting to you before it is
stated, because you have just spent an entire program on a quantity that
lacks it.

Program~\ref{prog:P30}'s divergence is emphatically \textbf{not} symmetric.
Its whole section~3 was about how much the two directions differ, and
section~5 about what choosing between them costs.

Is $I(X; Y)$ the same as $I(Y; X)$?
\dotline
\nextframe
\end{fr}

\begin{fr}
\ans{Yes, exactly}

And you can see it without computing anything. Swapping $X$ and $Y$ swaps
the two arguments of nothing: the joint $p(x, y)$ is the same table read the
other way, and the product $p(x)\,p(y)$ is the same product. Neither argument
of the divergence has moved.

Checked over the same $\val{p31.sym.trials}$ joints, with the two answers
agreeing to the last bit the arithmetic carries.

\begin{note}
It is worth being clear about what this does \emph{not} say. The divergence
is still asymmetric; it is the particular pair of arguments here that happens
to be unchanged under the swap. Section~3 shows the asymmetry surviving one
level down, in the pieces the quantity can be assembled from.
\end{note}
\end{fr}

% ============================================================
\section{Symmetric, out of asymmetric pieces}
\label{sec:P31-conditional}
\index{entropy!conditional}
\index{mutual information!as a reduction in uncertainty}

\begin{fr}
The definition is exact and it is not how anybody talks about the quantity.
People say \emph{how much knowing $X$ tells you about $Y$}, which sounds like
a different thing altogether.

One definition connects the two, and it is the only new object in this
program. The \textbf{conditional entropy}
\[ H(Y \mid X) \;=\; H(X, Y) \;-\; H(X) \]
is what is left of $Y$'s uncertainty once $X$ is known: the uncertainty in
the pair, less the uncertainty in $X$ alone.

Take Program~\ref{prog:P24}'s pair again, with $W = Z^{2}$. What is
$H(W \mid Z)$?
\dotline
\nextframe
\end{fr}

\begin{fr}
\ans{Exactly $0$}

$W$ is a function of $Z$. Once $Z$ is on the table there is no uncertainty
left in $W$ at all, and the arithmetic returns zero rather than something
small.

Now the same question the other way round, which is the one worth thinking
about: $H(Z \mid W)$. You are told $W$. How much of $Z$ is still open?
\dotline
\nextframe
\end{fr}

\begin{fr}
\ans{$\val{p31.zw.hzw}$ nats \dash{} emphatically not zero}

Told that $W = 1$, you know $Z$ is $-1$ or $1$ and you do not know which.
Told that $W = 0$, you know $Z$ exactly. On average
$\val{p31.zw.hzw}$ nats of $Z$ survive being told $W$.

So the two conditional entropies are as different as they can be: one is
exactly zero and the other is not.

\begin{trapbox}
\textbf{\enquote{Mutual information is symmetric, so the relationship is
symmetric} does not follow, and this pair is the counterexample.}

$Z$ determines $W$ and $W$ does not determine $Z$. That asymmetry is real,
it is what $H(W \mid Z) = 0$ against $H(Z \mid W) = \val{p31.zw.hzw}$ says,
and the mutual information cannot see it. Symmetry is a property of the
quantity, not of the relationship it is measuring.
\end{trapbox}
\end{fr}

\begin{fr}
Both of those conditional entropies reach the same mutual information, which
is the fact worth carrying out of this section:
\[ I(X; Y) \;=\; H(Y) - H(Y \mid X) \;=\; H(X) - H(X \mid Y) \]
\dash{} the uncertainty you had, less the uncertainty you are left with.
That is the sentence people mean when they say \emph{how much knowing $X$
tells you about $Y$}, and it is the same number as the divergence.

Check it on the pair, from the harder side. $H(Z)$ is $\val{p31.zw.hz}$ nats
and $H(Z \mid W)$ is $\val{p31.zw.hzw}$. What do they leave?
\dotline
\nextframe
\end{fr}

\begin{fr}
\ans{$\val{p31.zw.hz} - \val{p31.zw.hzw} = \val{p31.zw.mi}$ nats}

Which is what the other side gave: $H(W) - 0$, all of $W$'s uncertainty.
One number, reached from two sets of pieces that have almost nothing in
common.

\begin{aibox}
This is why the same quantity carries two names in two literatures and they
are the same thing. \emph{Information gain} in a decision tree is
$H(\text{label}) - H(\text{label} \mid \text{split})$, computed at every
candidate split and maximised. That is the right-hand form. \emph{Mutual
information} in a feature-selection routine is the divergence form. A
library that offers both is offering one number twice, and the two will
agree to whatever precision its estimator has \dash{} which is section~4's
subject and is less than people assume.
\end{aibox}

\mermaidfig{p31-two-routes-one-number}{Two ways to the same quantity}{one number}
\end{fr}

% ============================================================
\section{What a finite sample does to it}
\label{sec:P31-estimation}
\index{mutual information!estimation bias}
\index{estimation!of mutual information}

\begin{fr}
Everything so far has assumed you know the joint distribution. You do not.
What you have is a sample, and what you compute is the mutual information of
the table that sample makes \dash{} the \textbf{plug-in} estimate.

Program~\ref{prog:P26} spent a program on the gap between an estimator and
the thing it estimates, and on the two ways of being wrong. This is the same
question asked about this quantity.

Start with the easy case, and predict before reading on. Two quantities are
\emph{genuinely independent}, so the truth is exactly zero. You take a finite
sample and compute the plug-in estimate. What do you get?
\dotline
\nextframe
\end{fr}

\begin{fr}
\ans{A positive number, essentially always}

Not sometimes, and not with a small chance of being wrong in the other
direction: to get zero the sample table has to have its cells in \emph{exact}
proportion to its own margins, and a finite sample almost never does.

Anything else is a table with some dependence visible in it, and the plug-in
estimate reports that dependence, because that is what it was asked to do.

Here is the smallest instance of it. Twenty tosses of two fair coins, which
are independent by construction:

\transcript{p31-plug-in}

The dead-flat table reports zero. The one that is off by a single count in
each cell reports something else, and it is not small.
\end{fr}

\begin{fr}
That is one table. The question that matters is what the estimator does
\emph{on average}, because that is what a bias is.

For a two-by-two table with a fixed total there is a finite number of
possible outcomes \dash{} every way of splitting $N$ counts into four cells
\dash{} and each has a known probability. So the expectation can be computed
by adding up all of them, rather than by drawing samples and averaging.

At $N = 10$ there are $\val{p31.bias.tables10}$ such tables; at $N = 100$
there are $\val{p31.bias.tables100}$. Every one is enumerated and the
weights are checked to sum to exactly $1$.

Before the numbers: is this a measurement or a proof?
\dotline
\nextframe
\end{fr}

\begin{fr}
\ans{A proof, for that $N$}

Nothing was sampled and nothing was left out, so the answer is the
expectation and not an estimate of it. That is
Program~\ref{prog:P14}'s distinction, and it is the same discipline
Program~\ref{prog:P30} used on its own family of candidates: enumerate where
you can, so that the answer is not a statement about your sampler.

Here is what the enumeration gives, on data whose true mutual information is
exactly zero.

\begin{center}
\begin{tabular}{rrrr}
\toprule
$N$ & tables & $\Ex[\text{plug-in}]$ & vs $\frac{1}{2N}$ \\
\midrule
$10$ & $\val{p31.bias.tables10}$ & $\val{p31.bias.n10}$ & $\val{p31.bias.ratio10}$ \\
$20$ & $\val{p31.bias.tables20}$ & $\val{p31.bias.n20}$ & $\val{p31.bias.ratio20}$ \\
$50$ & $\val{p31.bias.tables50}$ & $\val{p31.bias.n50}$ & $\val{p31.bias.ratio50}$ \\
$100$ & $\val{p31.bias.tables100}$ & $\val{p31.bias.n100}$ & $\val{p31.bias.ratio100}$ \\
\bottomrule
\end{tabular}
\end{center}

Every entry in the third column is a positive number reported on data that
has nothing in it.
\end{fr}

\begin{fr}
The last column is the part worth slowing down for.

There is a standard correction for this bias. For a table with $a$ rows and
$b$ columns it subtracts $\frac{(a-1)(b-1)}{2N}$, which at two by two is
$\frac{1}{2N}$, and the column shows the true bias as a multiple of it.

Read down that column. What is the correction doing, and where is it doing
it worst?
\dotline
\nextframe
\end{fr}

\begin{fr}
\ans{Understating the bias, and worst at the smallest $N$}

$\val{p31.bias.ratio10}$ times the correction at $N = 10$, falling to
$\val{p31.bias.ratio100}$ by $N = 100$.

The correction is asymptotic: it is derived for large $N$ and it converges,
which the column shows it doing. The consequence is the one nobody states.

\begin{trapbox}
\textbf{\enquote{I applied the standard bias correction, so the number is
sound.}}

The correction is least accurate exactly where the bias is largest. At
$N = 100$ it is within two per cent and the bias it is removing is
$\val{p31.bias.n100}$ nats, which nobody needed removed. At $N = 10$ it
leaves nearly a third of the bias behind, and that is the regime the
correction was reached for.

A remedy that works where you do not need it and fails where you do is worse
than no remedy, because it removes the reason to be careful.
\end{trapbox}
\end{fr}

\begin{fr}
Now scale it, and scale it honestly.

That correction has an $(a-1)(b-1)$ in it, so the bias grows with the number
of \emph{cells} in the table and shrinks with the number of items. Sixteen
categories each way is $\val{p31.big.cells}$ cells.

The largest mutual information two sixteen-valued quantities can have is
$\ln \val{p31.big.k} = \val{p31.big.max}$ nats.

Here is the question a design review should ask. How many items do you need
before the bias is under one per cent of that maximum?
\dotline
\nextframe
\end{fr}

\begin{fr}
\ans{$\val{p31.big.need}$, which is about $\val{p31.big.per}$ items per cell}

At that sample size the bias is $\val{p31.big.bias}$ nats, one per cent of
the largest value the quantity can take. Below it, the number you compute is
part measurement and part sample size, in a proportion nothing on the page
tells you.

\begin{warning}
\textbf{This program quotes that figure the way round that is valid, and the
first draft did not.}

The obvious move is to run the formula at some small $N$ \dash{} sixteen
categories each way on fifty items predicts $\num{2.25}$ nats of bias, which
is most of the maximum, and it makes a striking sentence. It is also outside
the formula's own regime: fifty items over $\val{p31.big.cells}$ cells is a
fifth of an item per cell, and the formula is asymptotic in items against
cells. A probe of that case found it wrong by a third.

So the striking sentence is not available, and the sample size is, because
$\val{p31.big.need}$ items over $\val{p31.big.cells}$ cells is comfortably
inside the regime. Quoting a formula outside the range it was derived for is
the same error this section is about, one level up.
\end{warning}
\end{fr}

\begin{fr}
\begin{aibox}
Put those two numbers against a real setting. A probe reports the mutual
information between a layer's activations and a linguistic label, and the
evaluation set has a few hundred sentences.

The activations are continuous and get binned, or clustered, or quantised
\dash{} which fixes the number of cells, and nobody reports it. If that
number is anything like $\val{p31.big.cells}$ and the sample is a few
hundred, then by the arithmetic above the estimate is dominated by its own
bias, and it would have come out positive on shuffled labels.

There is a check that costs one extra run and settles it completely.
\end{aibox}

What is it?
\dotline
\nextframe
\end{fr}

\begin{fr}
\ans{Shuffle the labels and re-run the whole pipeline}

Whatever that reports is the floor: it is what your bins, your sample size
and your estimator produce on data with no relationship in it at all. The
claim is whatever is left above it.

It is almost never reported, and it is the difference between a number and a
result.

\mermaidfig{p31-bias-is-not-signal}{Where a plug-in estimate comes from}{signal or sample}
\end{fr}

% ============================================================
\section{Post-processing cannot create it}
\label{sec:P31-dpi}
\index{data-processing inequality}
\index{mutual information!data-processing inequality}

\begin{fr}
There is one theorem left and it is the one that does the most work.

Suppose $Z$ is computed from $Y$ and from nothing else \dash{} $Y$ is the
layer, $Z$ is what your probe reports, and the probe sees the layer and not
the input. Write that as $X \to Y \to Z$.

The probe may be anything: a linear classifier, a deep network, a lookup
table, a coin flip that ignores its input.

Take a guess before the frame turns over. Can $I(X; Z)$ exceed $I(X; Y)$?
\dotline
\nextframe
\end{fr}

\begin{fr}
\ans{No, never, whatever the probe is}

\[ I(X; Z) \;\le\; I(X; Y) \qquad \text{whenever } X \to Y \to Z \]

This is the \textbf{data-processing inequality}, and the informal reading is
the whole of it: \emph{you cannot create information by processing.} What
$Z$ knows about $X$, it knows through $Y$, so it cannot know more than $Y$
does.

Checked here rather than assumed. Against a fixed joint carrying
$\val{p31.dpi.ixy}$ nats, every channel from $Y$ to $Z$ on a rational grid
was enumerated \dash{} $\val{p31.dpi.channels}$ of them \dash{} and not one
exceeds it.
\end{fr}

\begin{fr}
That is a theorem about probes and it is worth reading slowly, because it
settles a claim people make constantly.

Somebody trains a better probe on the same layer and reports a higher score.
The layer has not changed and the data has not changed.

What has been learnt about the layer?
\dotline
\nextframe
\end{fr}

\begin{fr}
\ans{Nothing about the layer at all}

The inequality says every probe's score was bounded by $I(X; Y)$ the whole
time. A probe that scores higher has got closer to a ceiling that was always
there. It is a better probe, and the layer is what it was.

\begin{trapbox}
\textbf{\enquote{Our stronger probe recovers more syntax, so the
representation encodes syntax more richly than we thought.}}

The second clause does not follow from the first. The representation encodes
what it encodes; the probe got better at reading it. Those are different
sentences about different objects, and only one of them was measured.

The inequality is what makes this a theorem rather than a quibble: no amount
of probe capacity can push the score past $I(X; Y)$, so an improvement is
always movement towards the ceiling and never movement of it.
\end{trapbox}
\end{fr}

\begin{fr}
The same inequality has another end, and it is the one nobody quotes.

Go back to the $\val{p31.dpi.channels}$ enumerated probes. The best of them
is a real number, and it is not $\val{p31.dpi.ixy}$.

It reaches $\val{p31.dpi.best}$ nats. Against a layer carrying
$\val{p31.dpi.ixy}$, how much is the best available probe in that family
leaving behind?
\dotline
\nextframe
\end{fr}

\begin{fr}
\ans{More than a quarter of it}

And that is the \emph{best} of them. Every other probe in the family leaves
more.

So the inequality reads two ways and both are routinely got wrong, in
opposite directions, by the same people:

\begin{center}
\begin{tabular}{ll}
\toprule
\textbf{A better probe scores higher} & says nothing about the layer \\
\textbf{Every probe scores badly} & also says nothing about the layer \\
\bottomrule
\end{tabular}
\end{center}

A probe's score is a \textbf{lower bound} on what the layer carries. The
first row is the ceiling being approached; the second is the bound being
loose. Nothing in a probing result tells you which of the two you are looking
at, and the gap is not reported because it is not measured.
\end{fr}

\begin{fr}
It is worth being precise about why the gap exists here, because it is not a
failure of the probes.

The family enumerated above is every map from a three-valued $Y$ to a
two-valued $Z$.

Say what any such map is forced to do, and why that costs something.
\dotline
\nextframe
\end{fr}

\begin{fr}
\ans{It has to merge two of the three values, and whatever it merges, some
of what distinguished them is gone}

The best choice loses least; no choice loses nothing.

A probe with a smaller output than its input is in exactly that position, and
that describes every probe anybody trains.
\end{fr}

\begin{fr}
\begin{aibox}
The inequality also settles a question about layers, and this one goes the
way people do not expect.

A network computes layer $k+1$ from layer $k$. So
$X \to \text{layer } k \to \text{layer } k+1$, and the inequality applies:
\textbf{information about the input can only fall with depth, never rise.}

That sounds like it contradicts everything anybody says about deeper layers
carrying richer features, and it does not \dash{} it corrects what the claim
can mean. A deep layer cannot know more about the input. What it can be is
\emph{easier to read}, which is a statement about what a small probe can
extract, not about what is there. Those two are the two ends of this
section's inequality, and running them together is the mistake.
\end{aibox}

\mermaidfig{p31-probe-is-a-floor}{What a probe's score bounds}{a score is a floor}
\end{fr}

% ============================================================
\section{What would settle it}
\label{sec:P31-notmeasured}
\index{mutual information!what this book has not measured}

\begin{fr}
This program has been unusually hard on a class of claim, so it owes you a
statement of what it has and has not done.

What it has done is arithmetic: a definition, two identities, an exact
expectation over enumerated tables, and an inequality checked against every
channel in a family. None of it involves a trained model.

What it has \emph{not} done is measure a layer of anything.

\begin{warning}
\textbf{No number in this program was computed from a real model's
activations, and none of the claims about probing is a measurement.}

The arithmetic is exact and it constrains what a measurement could mean. It
does not tell you what any particular layer carries, and this book does not
have a trained model to ask. Programs \ref{prog:P08} and \ref{prog:P11} left
the same debt about rank and said so; this is a third entry in it.
\end{warning}
\end{fr}

\begin{fr}
So here is what a claim of the form \emph{layer 12 contains information about
syntax} would need in order to be supported, and every item is something this
program has given you a reason for.

Take a moment and list what you would ask for, before reading on.
\dotline
\nextframe
\end{fr}

\begin{fr}
\begin{ansblock}
\begin{enumerate}
\item \textbf{The number of cells}, which nobody reports \dash{} the binning
  or clustering that turned continuous activations into something countable.
  Section~4: the bias grows with cells and shrinks with items.
\item \textbf{The sample size against that number}, as items per cell.
\item \textbf{The shuffled-label floor}: the same pipeline run on labels with
  no relationship to the input. Whatever it reports is the bias, and the
  claim is only what sits above it.
\item \textbf{What the probe's capacity was}, because section~5 says a
  score is a lower bound and a bigger probe raises it without moving the
  layer.
\item \textbf{What is being compared with what.} A score alone is
  uninterpretable; a score against the shuffled floor, or against the same
  probe on another layer, is not.
\end{enumerate}
\end{ansblock}

Not one of the five is difficult. The third is the one that would settle
most disputes and it costs a single extra run.
\end{fr}

\begin{fr}
One last thing, and it is the honest counterweight.

None of this says the claims are false. Mutual information is the right
quantity for the question, it is the only one on the table that sees what a
correlation misses, and a layer that carries information about syntax is a
perfectly sensible thing for a layer to do.

What the program says is narrower and harder to argue with: \textbf{the usual
evidence does not distinguish the claim from its absence.} A positive
estimate is what you get either way, and the run that would tell them apart
is the one nobody does.

That is the difference between a claim being wrong and a claim being
unsupported, and only the second is being asserted here.
\end{fr}

\begin{summarybox}
\sumitem{1--6}{\textbf{A zero correlation settles nothing}:
  Program~\ref{prog:P24}'s pair has correlation exactly $0$ and mutual
  information $\val{p31.zw.mi}$ nats, which is all of $W$'s uncertainty.}
\sumitem{3--4}{\textbf{Mutual information is a divergence}:
  $I(X; Y) = \KL{p(x,y)}{p(x)p(y)}$, the gap between the joint you have and
  the product of margins you would have had under independence.}
\sumitem{7--8}{\textbf{So it is never negative and is zero exactly under
  independence}, both inherited from Program~\ref{prog:P30} with no new
  proof.}
\sumitem{9--10}{\textbf{It is symmetric}, because swapping the two
  quantities leaves both arguments of the divergence unchanged.}
\sumitem{11--15}{\textbf{The pieces it is built from are not}: on the same
  pair $H(W \mid Z) = 0$ and $H(Z \mid W) = \val{p31.zw.hzw}$, and both
  routes reach $\val{p31.zw.mi}$.}
\sumitem{16--19}{\textbf{The plug-in estimate is positive on independent
  data}, in expectation, at every sample size \dash{} $\val{p31.bias.n10}$
  nats at $N = 10$ by exact enumeration of all
  $\val{p31.bias.tables10}$ tables.}
\sumitem{20--21}{\textbf{The standard correction fails worst where it is
  needed most}: $\val{p31.bias.ratio10}$ times off at $N = 10$, converging to
  $\val{p31.bias.ratio100}$ by $N = 100$.}
\sumitem{22--23}{\textbf{Sixteen categories each way needs
  $\val{p31.big.need}$ items} for the bias to fall under one per cent of the
  maximum $\val{p31.big.max}$ nats \dash{} about $\val{p31.big.per}$ per
  cell.}
\sumitem{24--25}{\textbf{The check that settles it is one extra run}:
  shuffle the labels, re-run the whole pipeline, and report only what sits
  above whatever that returns.}
\sumitem{26--27}{\textbf{Post-processing cannot create information}: if
  $X \to Y \to Z$ then $I(X; Z) \le I(X; Y)$, checked against all
  $\val{p31.dpi.channels}$ channels in a family.}
\sumitem{28--33}{\textbf{So a probe's score is a lower bound}, and the best
  probe in that family still leaves more than a quarter behind. A better
  probe says nothing about the layer, and neither does a worse one.}
\sumitem{35--38}{\textbf{Nothing here was measured on a real model}, and the
  claim being made is that the usual evidence does not distinguish the claim
  from its absence \dash{} not that it is false.}
\end{summarybox}

\begin{testexercises}
\item $X$ is uniform on $\{-2, -1, 1, 2\}$ and $Y = \lvert X \rvert$.
  Without computing a divergence, say what $I(X; Y)$ is and why.
  \answerto{$H(Y) = \ln 2$ nats. $Y$ is a function of $X$, so
  $H(Y \mid X) = 0$ and $I(X; Y) = H(Y) - 0$. The same argument as
  Program~\ref{prog:P24}'s pair, with two values instead of three.}
\item Two quantities have $I(X; Y) = 0$. Somebody concludes their
  correlation is zero. Are they right, and is the converse also safe?
  \answerto{Right, and the converse is not. $I = 0$ is independence, which
  forces zero correlation. Zero correlation does not force $I = 0$, which is
  the whole of section~1.}
\item You compute a plug-in mutual information of $\num{0.05}$ nats between
  a binary feature and a binary label on $60$ items. What is the first thing
  to do?
  \answerto{Compare it against the bias on independent data at that size:
  the table in section~4 puts the expectation between $\val{p31.bias.n50}$
  and $\val{p31.bias.n100}$ nats there, so most of $\num{0.05}$ could be
  the estimator. Shuffle the labels, re-run, and report the difference.}
\item A colleague reports that a probe on layer 20 scores higher than the
  same probe on layer 4, and concludes layer 20 carries more information
  about the input. What is wrong with the conclusion?
  \answerto{Layer 20 is computed from layer 4, so the data-processing
  inequality says it cannot carry more information about the input \dash{}
  only less or the same. What changed is how easily a probe of that size can
  extract it.}
\item Write down $I(X; Y)$ three ways: as a divergence, and as a difference
  of entropies in each of the two directions.
  \answerto{$\KL{p(x,y)}{p(x)p(y)}$, and $H(Y) - H(Y \mid X)$, and
  $H(X) - H(X \mid Y)$. All three are the same number; the second and third
  are built from different pieces.}
\item Why can the plug-in estimator not be fixed by taking a larger sample
  alone, if the number of cells is allowed to grow with it?
  \answerto{The bias goes as cells over items, so growing both together
  leaves the ratio where it was. What has to fall is items per cell, which
  is why section~4 states the requirement that way round.}
\end{testexercises}

\canyou

\begin{furtherproblems}
\item Show that $I(X; X) = H(X)$, and say in words what that means.
  \answerto{$H(X \mid X) = 0$, so $I(X; X) = H(X) - 0 = H(X)$. A quantity
  tells you everything about itself, and \enquote{everything} is exactly its
  own entropy \dash{} which is why entropy is sometimes called
  self-information.}
\item Program~\ref{prog:P30} showed its divergence is not symmetric. Explain
  in one sentence why this one is, without computing anything.
  \answerto{Swapping $X$ and $Y$ leaves both arguments of the divergence
  unchanged: the joint is the same table and the product of the margins is
  the same product. The divergence is still asymmetric; this particular pair
  of arguments is fixed by the swap.}
\item A pipeline bins a continuous activation into $\val{p31.big.k}$ buckets
  and computes mutual information against a $\val{p31.big.k}$-way label.
  Somebody proposes doubling the buckets to \enquote{resolve more
  structure}. What does that do to the sample size needed?
  \answerto{Cells go up fourfold, so the bias at a fixed sample goes up
  about fourfold and the items needed to hold the bias fixed go up about
  fourfold too \dash{} from $\val{p31.big.need}$ to some sixteen thousand.
  Finer bins buy resolution and pay for it in items, and the payment is
  usually not made.}
\item The data-processing inequality says $I(X; Z) \le I(X; Y)$ when
  $X \to Y \to Z$. Give a $Z$ for which it holds with equality.
  \answerto{Any invertible function of $Y$ \dash{} relabelling the values,
  say. Nothing is merged, so nothing is lost. That is also the test for
  whether a preprocessing step can cost you anything: ask whether it can be
  undone.}
\item Section~4 measured the bias on \emph{independent} data. Argue that a
  plug-in estimate on \emph{dependent} data is still not trustworthy without
  the shuffled-label run.
  \answerto{The bias does not vanish when there is real dependence; it adds
  to it. So a positive estimate is a real quantity plus an unknown amount of
  estimator, and the shuffled run is what measures the second so it can be
  subtracted. Without it the two are not separable.}
\end{furtherproblems}
